\documentclass[12pt,a4paper]{article}

% ================= PACOTES =================
\usepackage[utf8]{inputenc}
\usepackage[T1]{fontenc}
\usepackage[brazil]{babel}
\usepackage{setspace}
\usepackage{geometry}
\usepackage{titlesec}
\usepackage{authblk}
\usepackage{indentfirst}
\usepackage{ragged2e}
\usepackage{abstract}
\usepackage{times}

% ================= MARGENS =================
\geometry{
  top=3cm,
  bottom=2cm,
  left=3cm,
  right=2cm
}

% ================= TÍTULOS =================
\titleformat{\section}{\bfseries}{\thesection.}{0.5em}{}
\titleformat{\subsection}{\bfseries}{\thesubsection}{0.5em}{}

% ================= DOCUMENTO =================
\begin{document}

% ================= TÍTULOS =================
\begin{center}
\textbf{TÍTULO EM PORTUGUÊS}\\[0.3cm]
\textbf{TITLE IN ENGLISH}\\[0.3cm]
\textbf{TÍTULO EN ESPAÑOL}
\end{center}

\vspace{0.5cm}

% ================= AUTORIA =================
\begin{flushright}
\textbf{Nome do Autor}\\
Titulação – Instituição – País\\
E-mail: autor@email.com
\end{flushright}

\vspace{0.5cm}

% ================= RESUMO =================
\begin{abstract}
\fontsize{10}{12}\selectfont
\justifying
Resumo em português, com no máximo 300 palavras, espaçamento simples, texto justificado e sem recuo.
\end{abstract}

\noindent\textbf{Palavras-chave:} Palavra 1; Palavra 2; Palavra 3.

\vspace{0.3cm}

% ================= ABSTRACT =================
\begin{abstract}
\fontsize{10}{12}\selectfont
\justifying
Abstract in English.
\end{abstract}

\noindent\textbf{Keywords:} Keyword 1; Keyword 2; Keyword 3.

\vspace{0.3cm}

% ================= RESUMEN =================
\begin{abstract}
\fontsize{10}{12}\selectfont
\justifying
Resumen en español.
\end{abstract}

\noindent\textbf{Palabras clave:} Palabra 1; Palabra 2; Palabra 3.

% ================= TEXTO =================
\section{Introdução}

Texto da introdução.

\subsection{Objetivos Gerais}

Texto dos objetivos.

\section{Revisão da Literatura}

Texto da revisão.

\section{Considerações Finais}

Texto das considerações finais.

% ================= REFERÊNCIAS =================
\section*{Referências}

\singlespacing
\noindent SOBRENOME, Nome. \textbf{Título do livro}. Local: Editora, ano.

\end{document}
